\documentclass[11pt,letterpaper]{article}

\usepackage[margin=0.85in,headheight=22pt]{geometry}
\usepackage[T1]{fontenc}
\usepackage[utf8]{inputenc}
\usepackage{lmodern}
\usepackage{microtype}
\usepackage{array}
\usepackage{booktabs}
\usepackage{tabularx}
\usepackage{longtable}
\usepackage{enumitem}
\usepackage{xcolor}
\usepackage{colortbl}
\usepackage{titlesec}
\usepackage{fancyhdr}
\usepackage[hidelinks]{hyperref}

\definecolor{inkblue}{HTML}{0B2545}
\definecolor{headingblue}{HTML}{2E74B5}
\definecolor{darkblue}{HTML}{1F4D78}
\definecolor{mutedgray}{HTML}{555555}
\definecolor{tablefill}{HTML}{F4F6F9}
\definecolor{rulegray}{HTML}{D0D7DE}

\setlength{\parindent}{0pt}
\setlength{\parskip}{7pt}
\renewcommand{\arraystretch}{1.25}
\setlist[itemize]{leftmargin=1.2em,itemsep=2pt,topsep=2pt}
\setlist[enumerate]{leftmargin=1.5em,itemsep=2pt,topsep=2pt}
\emergencystretch=2em

\titleformat{\section}{\large\bfseries\color{headingblue}}{\thesection}{0.7em}{}
\titleformat{\subsection}{\normalsize\bfseries\color{headingblue}}{\thesubsection}{0.7em}{}
\titleformat{\subsubsection}{\normalsize\bfseries\color{darkblue}}{\thesubsubsection}{0.7em}{}
\titlespacing*{\section}{0pt}{14pt}{5pt}
\titlespacing*{\subsection}{0pt}{10pt}{4pt}
\titlespacing*{\subsubsection}{0pt}{8pt}{3pt}

\pagestyle{fancy}
\fancyhf{}
\lhead{\footnotesize\color{mutedgray}Research Proposal}
\rhead{\footnotesize\color{mutedgray}Explainable VLMs for Construction Safety}
\cfoot{\footnotesize\thepage}
\renewcommand{\headrulewidth}{0.3pt}
\renewcommand{\headrule}{\hbox to\headwidth{\color{rulegray}\leaders\hrule height \headrulewidth\hfill}}

\newcolumntype{Y}{>{\raggedright\arraybackslash}X}
\newcommand{\tableheader}{\rowcolor{tablefill}}

\begin{document}

\begin{center}
{\Large\bfseries\color{inkblue} Research Proposal: Evaluating Explainability in Vision-Language Models for Dynamic Construction Safety Monitoring\par}
\vspace{5pt}
{\normalsize\color{mutedgray} Prepared for research discussion with Professor Mustafa Abdallah\par}
\end{center}

\vspace{5pt}
\noindent
\colorbox{tablefill}{%
  \parbox{\dimexpr\linewidth-2\fboxsep\relax}{%
  \textbf{\color{darkblue}Strategic focus:} Dynamic construction safety monitoring using ConstructionSite 10k, SODA, and CMA, with Florence-2 as the first pilot model.
  }%
}

\section{Introduction and Motivation}

Construction sites are complex and change quickly. Workers, equipment, materials, temporary structures, and environmental conditions interact in ways that make safety monitoring difficult. Manual inspections are important, but they are also time-consuming, subjective, and limited by what an inspector can observe at a given moment.

Traditional computer vision systems can detect objects such as hardhats, workers, excavators, ladders, or scaffolds. However, object detection alone does not explain whether a scene is safe. A model may detect a worker and a harness, but still fail to answer the safety question that matters: Is the worker tied off correctly while working at height?

Vision-Language Models (VLMs) offer a useful next step because they combine visual input with natural-language questions. They can support tasks such as:

\begin{itemize}
  \item answering safety questions about a site image,
  \item generating captions for observed hazards,
  \item grounding a safety answer to specific image regions,
  \item checking whether a scene appears compliant with a safety rule.
\end{itemize}

The problem is that VLMs can still make unreliable safety decisions. Their outputs can change when prompts are reworded, when lighting changes, when objects are partly hidden, or when the site is visually cluttered. In a safety setting, it is not enough for a model to say that a scene is safe or unsafe. A safety manager also needs to know what visual evidence the model used.

This project proposes an explainability evaluation framework for VLMs in construction safety. The framework adapts Professor Mustafa Abdallah's six-metric XAI evaluation approach from network intrusion detection to multimodal construction data. The goal is to test whether VLM explanations are accurate, focused, repeatable, efficient, resistant to noise, and valid across difficult safety cases.

Rather than claiming that every VLM can be fully opened and explained, this project will begin with open-source VLMs where internal visual representations can be inspected. The pilot will use gradient-based and attention-based attribution methods first, such as Integrated Gradients, Grad-CAM-style attribution, and attention rollout. More architecture-specific white-box methods, such as Layer-wise Relevance Propagation and DeepLIFT, will be added where the model structure makes them technically feasible.

\section{Data Foundation}

The proposed study will use a focused dataset stack that supports construction safety monitoring across spatial and temporal conditions.

\subsection{Dataset Stack}

\begingroup
\small
\renewcommand{\arraystretch}{1.15}
\begin{tabularx}{\linewidth}{>{\raggedright\arraybackslash}p{0.22\linewidth} >{\raggedright\arraybackslash}p{0.24\linewidth} Y}
\toprule
\tableheader \textbf{Dataset} & \textbf{Role in the Study} & \textbf{Why It Is Useful} \\
\midrule
ConstructionSite 10k & Primary safety VLM dataset & Provides construction-site images with safety-rule VQA, captions, grounding, bounding boxes, and image attributes. This is the strongest dataset for testing whether VLM answers are tied to the right visual evidence. \\
\midrule
SODA & Spatial clutter and object benchmark & Provides a large set of construction-site images with object annotations across varied layouts, weather, viewpoints, and site conditions. This supports object-level attribution and robustness testing. \\
\midrule
CMA, Construction Meta Action & Temporal safety behavior benchmark & Provides video clips of construction worker actions. This supports testing whether explanations remain useful when safety depends on motion or repeated behavior over time. \\
\bottomrule
\end{tabularx}
\endgroup

This stack keeps the project focused on dynamic safety monitoring. It avoids broadening the proposal into daily reporting or progress monitoring before the safety explainability question is well defined.

\subsection{Construction Anomaly Classes}

Professor Abdallah's work evaluates model explanations for normal network traffic and attack classes. In this project, the construction equivalent is safe site behavior versus safety anomalies.

\begingroup
\small
\renewcommand{\arraystretch}{1.12}
\begin{tabularx}{\linewidth}{>{\raggedright\arraybackslash}p{0.23\linewidth} >{\raggedright\arraybackslash}p{0.30\linewidth} Y}
\toprule
\tableheader \textbf{Class} & \textbf{Construction Meaning} & \textbf{Example} \\
\midrule
Compliant or safe state & No visible violation in the safety rule being checked & Worker has required PPE, equipment is separated from workers, access path is clear. \\
\midrule
PPE violation & Required protective equipment is missing or used incorrectly & No hardhat, no vest, harness visible but not tied off. \\
\midrule
Fall hazard & Worker is exposed to fall risk & Worker near unprotected edge, unsafe ladder use, missing guardrail. \\
\midrule
Struck-by or caught-between risk & Unsafe worker-equipment or worker-material interaction & Worker inside equipment swing radius, worker close to moving machine. \\
\midrule
Site obstruction or congestion & Unsafe site layout or blocked movement & Material pile blocks access, equipment crowds work area. \\
\midrule
Repeated unsafe action & Unsafe behavior persists over time & Repeated unsafe climbing, repeated close proximity to equipment, repeated missing PPE across frames. \\
\midrule
Subtle or occluded violation & Hazard is present but visually difficult & Worker partly hidden by dust, distance, glare, or overlapping equipment. \\
\bottomrule
\end{tabularx}
\endgroup

For the first pilot, the label set should stay small:

\begin{itemize}
  \item compliant or safe,
  \item PPE violation,
  \item fall hazard,
  \item struck-by or caught-between risk.
\end{itemize}

This limited scope is easier to test and easier to explain to the professor.

\section{Core Methodology}

The key translation is from tabular network features to multimodal construction features.

In Professor Abdallah's network security work, explanation methods rank features such as flow duration, destination port, packet length, or protocol type. In construction VLMs, the feature units are different:

\begin{itemize}
  \item visual regions, such as bounding boxes, image patches, object masks, or worker-equipment interaction zones,
  \item temporal features, such as object motion, action clips, or repeated unsafe states across frames,
  \item text features, such as safety-rule prompts, question tokens, and hazard-related words.
\end{itemize}

This project will treat the explanation unit as a multimodal feature pair. For example:

\begin{itemize}
  \item worker head region plus hardhat rule,
  \item harness region plus tie-off question,
  \item equipment swing area plus worker proximity rule,
  \item ladder region plus fall-protection prompt.
\end{itemize}

This gives the XAI framework a clear target. The explanation should show which visual regions and which rule concepts influenced the VLM's safety answer.

\section{Six-Metric XAI Evaluation Plan}

The study will evaluate VLM explanations using six metrics adapted from Professor Abdallah's framework.

\subsection{Descriptive Accuracy}

This metric tests whether the explanation identifies features that actually matter to the model.

\textbf{Process:}
\begin{itemize}
  \item Rank the most important visual regions and text tokens.
  \item Mask or blur the top-ranked visual regions.
  \item Remove or reword the top-ranked safety-rule tokens.
  \item Measure whether the VLM answer changes or confidence drops.
\end{itemize}

If removing the highlighted features changes the answer, the explanation is more likely to be faithful to the model's behavior.

\subsection{Sparsity}

This metric tests whether the explanation is focused.

A useful safety explanation should concentrate on the hazard region, not the whole image. For example, if the question is about hardhat compliance, the explanation should focus on the worker's head region, not on background scaffolding or unrelated equipment.

\subsection{Stability}

This metric tests whether the explanation is repeatable.

The same image and the same safety question will be submitted multiple times. The top visual regions and text tokens should remain similar across runs. Stability can be measured with overlap scores, Intersection over Union for highlighted regions, or rank correlation for feature lists.

\subsection{Efficiency}

This metric tests whether the explanation can be generated within a practical time budget.

The study will measure:

\begin{itemize}
  \item explanation runtime,
  \item memory use,
  \item number of forward and backward passes,
  \item throughput across a batch of images.
\end{itemize}

This matters because a future safety tool may need to run on site laptops, tablets, or edge devices.

\subsection{Robustness}

This metric tests whether explanations remain useful when the input is noisy or misleading.

The study will add controlled perturbations such as:

\begin{itemize}
  \item blur,
  \item low light,
  \item overexposure,
  \item partial occlusion,
  \item dust-like noise,
  \item viewpoint shift,
  \item misleading visual cues such as false or unclear PPE evidence.
\end{itemize}

The goal is to test whether the model explanation still points to the real safety cue, or whether the explanation is pulled toward irrelevant noise.

\subsection{Completeness}

This metric tests whether the explanation remains valid across all relevant sample types, including difficult cases.

The study will examine whether each sample receives a usable explanation and whether perturbing the top-ranked features changes the model's safety answer. If the model says that a region is important, but removing that region does not affect the answer, the explanation may be incomplete or unfaithful.

Completeness will be tested especially on:

\begin{itemize}
  \item crowded scenes,
  \item overlapping worker-equipment interactions,
  \item partial occlusions,
  \item subtle PPE misuse,
  \item multiple hazards in one image,
  \item video clips where the hazard depends on motion.
\end{itemize}

\section{Proposed 2 to 4 Week Pilot Experiment}

The pilot should be small and practical. It should show that the framework can run on real construction data without requiring full model training.

\subsection{Pilot Dataset}

Use 100 to 300 balanced samples from ConstructionSite 10k.

\textbf{Suggested sample groups:}
\begin{itemize}
  \item compliant scenes,
  \item PPE violations,
  \item fall hazards,
  \item struck-by or proximity risks.
\end{itemize}

\subsection{Pilot Model}

Use Florence-2 as the first pilot model.

\textbf{Reason:}
\begin{itemize}
  \item it is lightweight enough for a short pilot,
  \item it supports grounding and detection-style outputs,
  \item it is useful for connecting safety answers to image regions,
  \item it gives a practical starting point for visual attribution.
\end{itemize}

Qwen2-VL can be kept as a second-stage comparison model because it is stronger for general VQA and broader visual reasoning. However, Florence-2 is the better first model for a fast feasibility test.

\subsection{Pilot Steps}

\begin{enumerate}
  \item Select 100 to 300 samples from ConstructionSite 10k.
  \item Define a small set of safety-rule questions.
  \item Run Florence-2 on each image and prompt.
  \item Store the model answer, caption, grounding output, and confidence proxy if available.
  \item Generate visual attributions using gradient-based or attention-based methods.
  \item Mask top-ranked visual regions and re-run the model.
  \item Compute descriptive accuracy, sparsity, stability, and efficiency.
  \item Add a small robustness test using blur, lighting change, and occlusion.
  \item Compare highlighted regions against ground-truth boxes or expected safety regions.
\end{enumerate}

\section{Expected Contributions}

This project can make three clear contributions.

\begin{enumerate}
  \item A construction-specific XAI evaluation framework for VLM safety decisions.
  \item A dataset mapping that connects construction safety anomalies to Professor Abdallah's normal-versus-attack evaluation logic.
  \item A pilot benchmark showing whether VLM explanations remain accurate, focused, stable, efficient, and robust under construction-site conditions.
\end{enumerate}

The long-term deliverable is a journal paper that evaluates whether open VLMs can provide trustworthy explanations for safety-critical construction monitoring. A suitable target venue would be Automation in Construction, with possible alternatives in IEEE Access or related AI-for-construction venues.

\section{Collaboration Rationale}

This project fits both research areas.

\textbf{The construction-management side contributes:}
\begin{itemize}
  \item construction safety context,
  \item domain-specific hazard classes,
  \item construction datasets,
  \item VLM use cases,
  \item interpretation of site conditions.
\end{itemize}

\textbf{Professor Abdallah's side contributes:}
\begin{itemize}
  \item XAI evaluation methodology,
  \item experience with black-box and white-box explanation methods,
  \item robustness and completeness testing,
  \item statistical evaluation of explanation quality.
\end{itemize}

The proposed collaboration is therefore not just applying a VLM to construction images. It is testing whether the explanations behind VLM safety decisions can be trusted.

\end{document}
